\documentclass{article}
\title{ProjectGAN}
\author{Simeon Kost}
\date{Mar 2024}
\begin{document}
\maketitle
\section*{\centering\huge ProjectGAN}
\section*{What is ProjectGAN?}
\rule{\paperwidth}{0.4pt}
\\
\\ProjectGAN is all about making MLOps more efficient. 
It provides a complete infrastructure comprised of two layers. 
The first one gives developers a fully integrated UI on the Web with the ability to visualize analyze and operate a NN of various kinds making it possible to only concentrate on the development of your Neuronal Network. 
While the first layer is for data visualization and simple operations (feature engineering), the second layer is there to make the development of a NN architecture more efficient. 
This layer executes and supervises the underlaying program of your NN. 
Therefor it can handle errors, paus, interrupt, and rerun, … a NN with its corresponding code. 
You can interact with that layer through the terminal. 


\end{document}
